(1) 复束参数满足 \(1/q=1/R_c-i\lambda/(\pi w^2)\)。将 \(q=z+iz_R\) 取倒数，比较实虚部得 \(R_c(z)=z[1+(z_R/z)^2]\)（束腰处取无穷曲率半径），\(w(z)=w_0\sqrt{1+(z/z_R)^2}\)。基模的 Gouy 额外相位为 \(\psi(z)=\arctan(z/z_R)\)，从远前场至远后场变化约 \(\pi\)。

(2) 远离束腰 \(|z|\gg z_R\) 时光斑半径 \(w(z)\simeq w_0|z|/z_R\)，故远场发散半角 \(\theta\simeq w_0/z_R=\lambda/(\pi w_0)\)。于是 \(w_0\theta=\lambda/\pi\)，是理想高斯基模的衍射限定束腰—发散乘积。真实高阶或非理想激光光束用 \(M^2\ge1\) 修正，不能对任意非高斯模式直接套用理想基模极限。

(3) 在旧束腰处入射透镜有 \(q_{\rm in}=iz_R\)。ABCD 法则 \(q_{\rm out}=(Aq_{\rm in}+B)/(Cq_{\rm in}+D)=iz_R/(1-iz_R/f)\)。有理化得到 \(q_{\rm out}=[-z_R^2/f+iz_R]/[1+(z_R/f)^2]\)。出透镜后自由传播 \(q(z)=q_{\rm out}+z\)，因此新束腰位置为 \(z_f=z_R^2f/(f^2+z_R^2)>0\)，新瑞利长度为 \(z_R'=z_Rf^2/(f^2+z_R^2)\)。

(4) 当 \(f\gg z_R\)，有 \(z_f\simeq z_R^2/f\) 很小、\(z_R'\simeq z_R\)，透镜仅轻微改变束腰；当 \(f\ll z_R\)，有 \(z_f\simeq f\)、\(z_R'\simeq f^2/z_R\)，束腰压缩而发散增加。整个推导要求光束发散角及透镜引起的边缘光线倾角均远小于一弧度，且光学系统可近似为无像差薄透镜。
